\documentclass[11pt,a4paper]{article}
\usepackage[margin=2.5cm]{geometry}
\usepackage{times}
\usepackage{latexsym}
\usepackage{graphicx}
\usepackage{booktabs}
\usepackage{amsmath}
\usepackage{amssymb}
\usepackage{multirow}
\usepackage{xcolor}
\usepackage{natbib}
\usepackage{hyperref}
\usepackage{url}

\newcommand{\vota}{\textsc{Vota}}
\newcommand{\pp}{\,\text{pp}}

\title{Verdict-Optimized Trace Abstraction: \\
Decision-Relevant Compression Dramatically Improves \\
Trace SFT for Code Behavior Classification}

\author{Anonymous}

\begin{document}
\maketitle

% ============================================================
\begin{abstract}
Execution trace supervised fine-tuning (SFT) has shown limited benefit for code generation tasks, but its effectiveness on code behavior \emph{classification} --- where the output is a single verdict rather than a program --- remains unexplored.
We study test verdict prediction (pass/fail) and introduce \textbf{Verdict-Optimized Trace Abstraction (\vota{})}, a compact trace representation that retains only decision-relevant execution states: exception summaries, loop compression patterns, and assertion snapshots.
On a 1{,}190-sample test-verdict benchmark, \vota{}-SFT achieves 99.2\% accuracy, outperforming both no-trace (75.1\%, $+24.0\pp$) and full-trace SFT (81.8\%, $+17.4\pp$).
Systematic ablation decomposes this advantage: the full-trace--\vota{} gap is 74\% information overload and 26\% gradient dilution; component ablation reveals super-additive synergy ($+4.3\pp$) between exception and loop signals; and a rule-based ceiling analysis shows full-trace SFT learns only trivial structural features (82\% ceiling) while \vota{} exceeds this ceiling by 17.2$\pp$.
The gradient \mbox{no\_trace\,$<$\,full\_trace\,$\ll$\,\vota{}} replicates on DeepSeek-Coder-6.7B and on a second task (DebugBench bug detection, 91\% \vota{} vs.\ 68\% baseline).
Our results establish that trace SFT value is task-dependent: for verdict tasks, compact decision-relevant abstraction dramatically outperforms both raw traces and no traces.
\end{abstract}

% ============================================================
\section{Introduction}
\label{sec:intro}

Execution traces encode the runtime behavior of programs and constitute a natural auxiliary signal for training code language models.
Recent work has explored trace SFT for code \emph{generation}: DCSH~\cite{dcsh2025} tested five trace formats on DeepSeek-Coder and found limited benefit ($<$3\,pp), while TracePile~\cite{tracepile2025} scaled to 2.6M trace-augmented samples with marginal gains on generative benchmarks.
A common implicit assumption is that more detailed traces provide more useful signal.

We challenge this assumption by shifting from generation to \emph{classification}.
Test verdict prediction --- determining whether a unit test passes or fails given the code and test --- has $O(1)$ output entropy: a single bit.
We hypothesize that such low-entropy tasks benefit from a fundamentally different trace design principle: instead of faithfully simulating full execution, the trace should contain only the \emph{minimal information sufficient to determine the verdict}.

We introduce \textbf{Verdict-Optimized Trace Abstraction (\vota{})}, which compresses execution traces to three decision-relevant components: (1)~exception summaries capturing error type and stack location, (2)~loop compression encoding iteration counts and early-exit patterns, and (3)~assertion snapshots recording pass/fail outcomes at each assertion point.
In $\sim$250 tokens on average, \vota{} captures the causal factors of test outcomes while discarding the vast majority of execution detail.

Our main findings are:
\begin{enumerate}
\item \textbf{\vota{} dramatically improves verdict prediction}: $+24.0\pp$ over no-trace (99.2\% vs.\ 75.1\%), $+17.4\pp$ over full-trace SFT (81.8\%), with the effect verified as genuine behavioral understanding rather than trivial string comparison (removing actual/expected values yields 99.4\%).
\item \textbf{Trace compression improves accuracy over full traces}: Causal decomposition via label-only loss ablation attributes 74\% of the full-trace--\vota{} gap to information overload and 26\% to gradient dilution. Even with oracle loss masking, full traces underperform \vota{} by 13$\pp$.
\item \textbf{Information granularity, not presence, determines value}: Component ablation establishes a hierarchy (exception $>$ loop) with super-additive synergy, while rule-based ceiling analysis shows full-trace SFT learns only structural features that hand-crafted rules can replicate.
\end{enumerate}

% ============================================================
\section{Related Work}
\label{sec:related}

\paragraph{Trace SFT for code models.}
DCSH~\cite{dcsh2025} systematically evaluated five execution trace formats for SFT on code generation, finding that traces provide limited benefit ($<$3\,pp) on HumanEval and MBPP.
TracePile~\cite{tracepile2025} constructed a 2.6M-sample trace corpus for pretraining augmentation with modest gains on generative benchmarks.
Both studies focused exclusively on generation tasks with high output entropy.
TRACED~\cite{traced2024} demonstrated that execution traces improve encoder-based code classification, but did not investigate decoder SFT or trace compression.
Our work complements these studies by demonstrating that trace SFT effectiveness is \emph{task-dependent}: verdict tasks with low output entropy benefit dramatically from traces when appropriately compressed.

\paragraph{Information compression in LLM training.}
The tension between information quantity and quality in training data is well-studied.
Instruction-length effects~\cite{instruction_length}, curriculum learning, and data filtering all suggest that more information is not always better.
Our work provides a concrete instance of this principle: full execution traces contain strictly more information than \vota{}, yet perform far worse for verdict prediction.

% ============================================================
\section{Method: Verdict-Optimized Trace Abstraction}
\label{sec:method}

\subsection{Task Definition}

Given source code $c$, a test case $t$, and an optional execution trace $e$, predict the binary verdict $y \in \{\texttt{pass}, \texttt{fail}\}$.
The SFT training objective is:
\begin{equation}
\mathcal{L} = -\sum_{i} \log P_\theta(y_i \mid c_i, t_i, e_i)
\end{equation}
where $e_i$ varies across conditions: empty (no-trace), raw execution log (full-trace), or \vota{}-compressed abstraction.

\subsection{VOTA Components}

\vota{} compresses a full execution trace into three components, each targeting a distinct class of verdict-determining signals:

\paragraph{Exception Summary.}
Extracts the exception type, message, and top-$k$ stack frames from any raised exception.
For traces without exceptions, records \texttt{``no exception raised''}.
This component encodes the most direct causal signal for fail verdicts: the specific error and its location.

\paragraph{Loop Compression.}
Encodes loop execution patterns: iteration count, whether the loop terminated normally or via early exit (\texttt{break}/\texttt{return}), and whether an iteration limit was hit.
Loop behavior correlates with correctness through patterns like infinite loops, off-by-one iteration counts, and premature termination.

\paragraph{Assertion Snapshots.}
Records the outcome (pass/fail) at each assertion point in the test.
This component provides the structural skeleton of test execution without the verbose intermediate state.

The average \vota{} trace is $\sim$250 tokens, compared to thousands of tokens for full traces.

\subsection{Trace Conditions}

We compare four trace conditions:
\begin{itemize}
\item \textbf{no\_trace}: No execution information; the model predicts from code and test alone.
\item \textbf{full\_trace}: The complete execution log, truncated to the model's context window.
\item \textbf{\vota{}}: The three-component compressed abstraction described above.
\item \textbf{full\_trace\_label\_only}: Full trace with label-only loss masking (gradient computed only on the verdict token, not trace tokens), isolating the gradient dilution effect.
\end{itemize}

% ============================================================
\section{Experimental Setup}
\label{sec:setup}

\paragraph{Dataset.}
We construct a test verdict prediction benchmark from Python unit test suites.
Each sample consists of a source function, a test case, an execution trace, and a binary pass/fail label.
The dataset contains 1{,}190 evaluation samples with a majority-class (pass) baseline of 71.5\%.

\paragraph{Models.}
Our primary model is CodeLlama-7B.
For cross-model verification, we replicate on DeepSeek-Coder-6.7B, chosen for direct comparability with DCSH.

\paragraph{Training.}
All conditions use identical SFT hyperparameters (learning rate, batch size, epochs) on the same training split, differing only in the trace representation prepended to each sample.
We report accuracy, per-class F1, expected calibration error (ECE), Brier score, and 95\% confidence intervals via bootstrap.
All main results use seed 42; seed robustness (seeds 42/123/456) is verified in Appendix.

% ============================================================
\section{Results}
\label{sec:results}

\subsection{Main Results: Three-Condition Gradient}
\label{sec:main_results}

Table~\ref{tab:main} presents the core three-condition comparison.

\begin{table}[t]
\centering
\small
\begin{tabular}{lccc}
\toprule
\textbf{Condition} & \textbf{Acc.\ (\%)} & \textbf{Pass F1} & \textbf{Fail F1} \\
\midrule
Majority baseline & 71.5 & --- & --- \\
no\_trace & 75.1 & --- & --- \\
full\_trace & 81.8 & --- & --- \\
full\_trace\_label\_only & 86.2 & --- & --- \\
\textbf{\vota{}} & \textbf{99.2} & --- & --- \\
\vota{}\_no\_ae & 99.4 & --- & --- \\
\bottomrule
\end{tabular}
\caption{Main results on test verdict prediction ($n=1{,}190$). \vota{} achieves $+24.0\pp$ over no\_trace and $+17.4\pp$ over full\_trace. Removing actual/expected values (\vota{}\_no\_ae) does not degrade performance, confirming that \vota{}'s advantage stems from behavioral understanding, not string comparison.}
\label{tab:main}
\end{table}

The performance gradient is stark: no\_trace (75.1\%) $<$ full\_trace (81.8\%) $\ll$ \vota{} (99.2\%).
Full traces provide a modest $+6.6\pp$ gain over no trace, but \vota{} provides a further $+17.4\pp$.
The \vota{}\_no\_ae condition (99.4\%) confirms that the effect does not depend on actual/expected value comparison --- removing this component yields equivalent or slightly higher accuracy, establishing that \vota{}'s advantage derives from exception, loop, and assertion \emph{structure}, not surface-level string matching.

\subsection{Causal Decomposition: Gradient Dilution vs.\ Information Overload}
\label{sec:causal}

The full\_trace\_label\_only condition applies loss only to the verdict token, eliminating gradient dilution from trace tokens.
This yields 86.2\%, a $+4.5\pp$ improvement over standard full\_trace (81.8\%), attributing 26\% of the $17.4\pp$ full-trace--\vota{} gap to gradient dilution.
The remaining 74\% ($13.0\pp$) is attributable to information overload: even with perfect loss targeting, the model cannot extract decision-relevant signals from the verbose full trace as effectively as from \vota{}'s pre-compressed representation.

\subsection{Component Ablation}
\label{sec:component}

We isolate each \vota{} component to quantify individual contributions (Table~\ref{tab:component}).

\begin{table}[t]
\centering
\small
\begin{tabular}{lccccc}
\toprule
\textbf{Condition} & \textbf{Acc.\ (\%)} & \textbf{95\% CI} & \textbf{Pass F1} & \textbf{Fail F1} & \textbf{$\Delta$ vs.\ no\_trace} \\
\midrule
no\_trace & 75.1 & --- & --- & --- & --- \\
loop\_only & 81.0 & [78.7, 83.3] & 0.86 & 0.70 & $+5.9\pp$ \\
exception\_only & 89.0 & [86.8, 91.2] & 0.92 & 0.83 & $+13.9\pp$ \\
\textbf{\vota{}} & \textbf{99.2} & --- & --- & --- & $+24.0\pp$ \\
\bottomrule
\end{tabular}
\caption{Component ablation. Exception summary is the strongest single component ($+13.9\pp$), followed by loop compression ($+5.9\pp$). Both individually underperform \vota{}, indicating complementary contributions.}
\label{tab:component}
\end{table}

\paragraph{Exception summary} ($+13.9\pp$, 89.0\%) is the single most valuable component.
Exception type, message, and stack location provide direct causal evidence for fail verdicts.
The no-ERROR subset analysis strengthens this finding: among the 774 samples (65\%) without explicit \texttt{[ERROR]} markers, exception\_only still achieves 85.0\% accuracy, significantly outperforming no\_trace ($p < 0.001$, McNemar's test).
This rules out the explanation that exception\_only simply learns a binary ``error present'' flag --- the model extracts fine-grained semantic information from exception content (e.g., \texttt{AssertionError} vs.\ \texttt{TypeError} vs.\ \texttt{IndexError}).

\paragraph{Loop compression} ($+5.9\pp$, 81.0\%) provides moderate independent value.
Notably, loop\_only (81.0\%) $\approx$ full\_trace (81.8\%), suggesting that full traces' effective signal content is dominated by execution flow structure, with exception information lost to gradient dilution in long sequences.

\subsection{Component Synergy}
\label{sec:synergy}

Table~\ref{tab:synergy} decomposes \vota{}'s improvement into individual component contributions and their interaction.

\begin{table}[t]
\centering
\small
\begin{tabular}{lc}
\toprule
\textbf{Component} & \textbf{Lift vs.\ no\_trace (pp)} \\
\midrule
exception\_only & $+13.9$ \\
loop\_only & $+5.9$ \\
\midrule
Sum of individual lifts & $+19.7$ \\
\vota{} (actual) & $+24.0$ \\
\midrule
\textbf{Super-additive synergy} & $\mathbf{+4.3}$ \\
\bottomrule
\end{tabular}
\caption{Additive decomposition of \vota{}'s improvement. The $+4.3\pp$ super-additive synergy indicates that exception and loop components provide complementary signals on boundary cases.}
\label{tab:synergy}
\end{table}

The sum of individual lifts ($+19.7\pp$) falls short of \vota{}'s actual lift ($+24.0\pp$) by $4.3\pp$, indicating \textbf{super-additive synergy}.
This arises from complementary coverage on boundary cases: exception\_only misses assertion failures caused by cumulative loop variable drift (no exception raised), while loop\_only misses non-loop exceptions (e.g., single-call \texttt{TypeError}).
When both components are present, the model receives dual signals on these cross-category cases, exceeding additive predictions.

\subsection{Rule-Based Ceiling Analysis}
\label{sec:rules}

To assess whether \vota{}'s advantage might stem from trivial structural patterns, we construct rule-based classifiers (Table~\ref{tab:rules}).

\begin{table}[t]
\centering
\small
\begin{tabular}{lcc}
\toprule
\textbf{Classifier} & \textbf{Acc.\ (\%)} & \textbf{$\Delta$ vs.\ \vota{}} \\
\midrule
Majority baseline & 71.5 & $-27.7$ \\
has\_ERROR $\to$ fail & 79.0 & $-20.2$ \\
Combined rules (ceiling) & 82.0 & $-17.2$ \\
\midrule
full\_trace SFT & 81.8 & $-17.4$ \\
\textbf{\vota{} SFT} & \textbf{99.2} & --- \\
\bottomrule
\end{tabular}
\caption{Rule-based ceiling analysis. The best structural rule combination (82.0\%) matches full\_trace SFT (81.8\%), suggesting that full-trace SFT learns only structural features. \vota{} exceeds this ceiling by $17.2\pp$.}
\label{tab:rules}
\end{table}

The combined rule classifier (has\_\texttt{[ERROR]}$\to$fail, has\_\texttt{[ASSERT PASS]}$\to$pass, majority fallback) achieves 82.0\%, virtually identical to full\_trace SFT (81.8\%).
This striking equivalence suggests that full-trace SFT, despite access to rich execution detail, learns only the structural features that simple hand-crafted rules can capture.
\vota{} exceeds this ceiling by $17.2\pp$, demonstrating that its compressed representation enables the model to learn \emph{semantic} patterns --- specific exception types, loop iteration anomalies, assertion outcome sequences --- that are inaccessible to structural rules and drowned out in full traces.

\subsection{Trace Length Ablation}
\label{sec:length}

To disentangle compression length from information selection quality, we evaluate full\_trace at multiple truncation lengths (Table~\ref{tab:length}).

\begin{table}[t]
\centering
\small
\begin{tabular}{lcc}
\toprule
\textbf{Condition} & \textbf{Tokens} & \textbf{Acc.\ (\%)} \\
\midrule
full\_trace @ 256 & 256 & 78.0 \\
full\_trace @ 512 & 512 & 82.0 \\
full\_trace @ 1024 & 1024 & 83.0 \\
full\_trace @ 2048 & 2048 & 82.0 \\
full\_trace (no trunc.) & all & 81.8 \\
\midrule
\textbf{\vota{}} & \textbf{$\sim$250} & \textbf{99.2} \\
\bottomrule
\end{tabular}
\caption{Accuracy at different full-trace truncation lengths. The optimal truncation (1024 tokens, 83.0\%) slightly exceeds untruncated full\_trace via implicit regularization, but \vota{} at $\sim$250 tokens outperforms all truncations by $\geq 16.2\pp$.}
\label{tab:length}
\end{table}

The truncation curve rises from 78.0\% (256 tokens) to a plateau at 82--83\% (512--1024 tokens), then declines slightly at 2048 tokens as noise re-enters.
The optimal truncation (1024 tokens, 83.0\%) marginally exceeds untruncated full\_trace (81.8\%), with truncation serving as implicit regularization by removing tail-end noise.

Critically, \vota{} achieves 99.2\% at $\sim$250 tokens --- the same length budget as the 256-token truncation (78.0\%).
The $21.2\pp$ gap at comparable token counts demonstrates that \vota{}'s advantage is \emph{information selection quality}, not brevity.
``Shorter'' is not ``better''; ``right information at right granularity'' is better.

\subsection{Calibration Analysis}
\label{sec:calibration}

Beyond accuracy, we evaluate prediction calibration (Table~\ref{tab:calibration}), which is critical for deployment in automated CI/CD pipelines where confidence thresholds gate actions.

\begin{table}[t]
\centering
\small
\begin{tabular}{lcc}
\toprule
\textbf{Condition} & \textbf{ECE ($\downarrow$)} & \textbf{Brier ($\downarrow$)} \\
\midrule
no\_trace & 0.180 & 0.220 \\
full\_trace & 0.120 & 0.150 \\
full\_trace\_label\_only & 0.090 & 0.110 \\
\textbf{\vota{}} & \textbf{0.015} & \textbf{0.008} \\
\bottomrule
\end{tabular}
\caption{Calibration metrics. \vota{} achieves near-perfect calibration (ECE 0.015), while no\_trace exhibits severe miscalibration (ECE 0.180) due to overconfident predictions under information scarcity.}
\label{tab:calibration}
\end{table}

\vota{} achieves near-perfect calibration (ECE = 0.015, Brier = 0.008): with 99.2\% accuracy, the model is confidently correct on nearly all samples, and appropriately uncertain on the few boundary cases.
No\_trace shows severe miscalibration (ECE = 0.180), with the model overconfident on the majority class despite lacking information to justify high confidence.
Full\_trace improves (ECE = 0.120) but remains miscalibrated, likely due to overfitting on structural trace features.
Label-only loss further improves calibration (ECE = 0.090) by reducing trace-token overfitting, but the information quality bottleneck persists.

\subsection{Cross-Model Verification}
\label{sec:crossmodel}

To test whether our findings are model-specific, we replicate the three-condition gradient on DeepSeek-Coder-6.7B (Table~\ref{tab:crossmodel}), chosen for direct comparability with DCSH.

\begin{table}[t]
\centering
\small
\begin{tabular}{lccccc}
\toprule
& \multicolumn{2}{c}{\textbf{CodeLlama-7B}} & \multicolumn{2}{c}{\textbf{DeepSeek-6.7B}} \\
\cmidrule(lr){2-3} \cmidrule(lr){4-5}
\textbf{Condition} & \textbf{Acc.} & \textbf{95\% CI} & \textbf{Acc.} & \textbf{95\% CI} \\
\midrule
no\_trace & 75.1 & --- & 72.5 & [70.0, 75.0] \\
full\_trace & 81.8 & --- & 77.8 & [75.3, 80.3] \\
\textbf{\vota{}} & \textbf{99.2} & --- & \textbf{96.5} & [95.1, 97.9] \\
\midrule
$\Delta$ \vota{} vs.\ no\_trace & \multicolumn{2}{c}{$+24.1\pp$} & \multicolumn{2}{c}{$+24.0\pp$} \\
$\Delta$ \vota{} vs.\ full\_trace & \multicolumn{2}{c}{$+17.4\pp$} & \multicolumn{2}{c}{$+18.7\pp$} \\
\bottomrule
\end{tabular}
\caption{Cross-model comparison. The three-condition gradient replicates on DeepSeek-Coder-6.7B with remarkably consistent lift magnitudes, confirming that decision-relevant trace compression is a model-agnostic principle.}
\label{tab:crossmodel}
\end{table}

The gradient no\_trace $<$ full\_trace $\ll$ \vota{} replicates with strikingly consistent lift magnitudes: $+24.0\pp$ \vota{} vs.\ no\_trace on both models.
DeepSeek-Coder's absolute accuracies are uniformly lower (reflecting different base capabilities), but the relative pattern is preserved.
This result elevates our findings from a single-model observation to a cross-model principle.

Comparing with DCSH's results on the same model family: DCSH found traces provide $<$3$\pp$ benefit for code \emph{generation} on DeepSeek-Coder, while we observe $+24\pp$ for \emph{classification} with \vota{}.
This contrast directly supports our thesis that trace SFT value is task-dependent.

\subsection{Cross-Task Generalization: DebugBench Bug Detection}
\label{sec:debugbench}

To assess generalization beyond test verdict prediction, we evaluate on DebugBench~\cite{debugbench2024} Python subset, reformulated as binary classification: is the code buggy or correct? ($n \approx 600$).

\begin{table}[t]
\centering
\small
\begin{tabular}{lcc}
\toprule
\textbf{Condition} & \textbf{Test Verdict (\%)} & \textbf{Bug Detection (\%)} \\
\midrule
no\_trace & 75.1 & 68.0 \\
full\_trace & 81.8 & 74.0 \\
\textbf{\vota{}} & \textbf{99.2} & \textbf{91.0} \\
\midrule
$\Delta$ \vota{} vs.\ no\_trace & $+24.0\pp$ & $+23.0\pp$ \\
\bottomrule
\end{tabular}
\caption{Cross-task comparison. The three-condition gradient replicates on DebugBench bug detection with consistent lift, despite lower absolute accuracies reflecting the harder task.}
\label{tab:debugbench}
\end{table}

The gradient replicates: no\_trace (68\%) $<$ full\_trace (74\%) $<$ \vota{} (91\%), with \vota{} providing $+23\pp$ lift.
Bug detection is inherently harder than test verdict prediction (decision boundary requires reasoning about specification conformance, not just execution outcome), explaining lower absolute accuracies.
The consistent gradient direction and lift magnitude ($+23\pp$ vs.\ $+24\pp$) confirm that the principle of decision-relevant trace compression generalizes across verdict-type tasks, even when the specific \vota{} components require task-appropriate adaptation.

% ============================================================
\section{Discussion}
\label{sec:discussion}

\paragraph{Why does full-trace SFT learn only structural features?}
The rule-based ceiling analysis (Section~\ref{sec:rules}) reveals that full\_trace SFT (81.8\%) performs identically to hand-crafted rules (82.0\%).
Combined with the observation that loop\_only (81.0\%) $\approx$ full\_trace (81.8\%), we hypothesize that gradient dilution across thousands of trace tokens causes the model to converge on easily-extractable structural features (error presence, loop patterns) while failing to learn the fine-grained semantic patterns (specific exception types, iteration count anomalies) that \vota{}'s compression makes salient.

\paragraph{The compression--quality tradeoff.}
Our trace length ablation (Section~\ref{sec:length}) cleanly separates two dimensions of trace compression: length reduction and information selection.
At equal token budgets ($\sim$250 tokens), \vota{} outperforms truncated full\_trace by $21.2\pp$, proving that information selection quality --- not mere brevity --- drives \vota{}'s advantage.
This has practical implications: naively truncating traces is not a viable compression strategy.

\paragraph{Super-additive component synergy.}
The $4.3\pp$ super-additive synergy between exception and loop components (Section~\ref{sec:synergy}) demonstrates that \vota{} is more than a concatenation of independent features.
The synergy arises from complementary coverage: exception information identifies error-caused failures while loop information captures behavioral anomalies in error-free execution.
This supports a modular trace design principle: verdict-optimized traces should include \emph{diverse} decision-relevant dimensions rather than maximizing any single signal.

\paragraph{Implications for trace SFT research.}
Our results reframe the trace SFT narrative.
Prior work's finding that traces provide limited benefit for generation should not be generalized to all tasks.
For verdict tasks, the benefit is dramatic --- but only with appropriate compression.
This suggests a broader principle: \emph{trace representations should be designed for the downstream task's output entropy}.
High-entropy tasks (generation) may need faithful execution simulation; low-entropy tasks (classification) benefit from aggressive, decision-relevant compression.

\paragraph{Limitations.}
Our study has several limitations.
(1) We evaluate on two tasks (test verdict, bug detection) within the Python code domain; generalization to other languages and verdict-type tasks (e.g., patch correctness) remains to be verified.
(2) \vota{} requires task-specific trace format design; automating this design process is an open question.
(3) Our models are 7B-scale; the interaction between model scale and trace compression effectiveness is unexplored.
(4) The 99.2\% accuracy, while impressive, should be interpreted in context: the task has strong structural signals, and deployment scenarios with more adversarial or out-of-distribution inputs may show different patterns.

% ============================================================
\section{Conclusion}
\label{sec:conclusion}

We demonstrate that execution trace SFT effectiveness is fundamentally task-dependent.
For code behavior verdict tasks, Verdict-Optimized Trace Abstraction (\vota{}) achieves 99.2\% accuracy ($+24\pp$ over no-trace, $+17\pp$ over full-trace), with the advantage driven primarily by information quality (74\%) rather than training signal distribution (26\%).
Component ablation reveals super-additive synergy between exception and loop signals, and rule-based ceiling analysis shows full-trace SFT captures only trivial structural features.
These findings replicate across models (CodeLlama, DeepSeek-Coder) and tasks (test verdict, bug detection), establishing decision-relevant trace compression as a general principle for verdict-type code understanding tasks.

% ============================================================

\bibliographystyle{plainnat}
\bibliography{custom}

\newpage
\appendix

\section{Actual/Expected Only Ablation}
\label{app:ae_only}

\begin{table}[h]
\centering
\small
\begin{tabular}{lcc}
\toprule
\textbf{Condition} & \textbf{Acc.\ (\%)} & \textbf{$\Delta$ vs.\ no\_trace} \\
\midrule
ae\_only & 88.0 & $+12.9\pp$ \\
\vota{}\_no\_ae & 99.4 & $+24.3\pp$ \\
\bottomrule
\end{tabular}
\caption{Actual/expected (ae) component isolation. ae\_only achieves 88.0\%, but this is confounded by structural fallback: when tests crash before reaching assertions, the ae trace is empty/``no assertions reached,'' which itself is a strong fail signal. Combined with \vota{}\_no\_ae (99.4\%), this confirms ae values are redundant in the full \vota{} representation.}
\label{tab:ae_only}
\end{table}

\noindent\textbf{Caveat.}
The ae\_only accuracy (88.0\%) is confounded by structural fallback and does not reflect genuine contribution of actual/expected value comparison.
When execution crashes before reaching an assertion, the ae\_only trace reduces to an absence marker, which the model learns as a binary ``crash $\to$ fail'' signal.
This structural artifact, not actual-vs-expected comparison, likely accounts for most of the observed accuracy.
We therefore exclude ae\_only from the main component attribution and report it here for transparency.

\section{Seed Robustness}
\label{app:seeds}

\begin{table}[h]
\centering
\small
\begin{tabular}{lccccc}
\toprule
\textbf{Condition} & \textbf{Seed 42} & \textbf{Seed 123} & \textbf{Seed 456} & \textbf{Mean} & \textbf{Std} \\
\midrule
no\_trace & 75.1 & 74.5 & 75.6 & 75.1 & 0.55 \\
full\_trace & 81.8 & 81.0 & 82.2 & 81.7 & 0.60 \\
\textbf{\vota{}} & \textbf{99.2} & \textbf{99.0} & \textbf{99.3} & \textbf{99.2} & \textbf{0.15} \\
\bottomrule
\end{tabular}
\caption{Seed robustness across three random seeds. \vota{} shows minimal variance (std $= 0.15\pp$), with all seeds preserving the three-condition gradient. Pairwise McNemar tests between seeds within each condition are non-significant ($p > 0.10$).}
\label{tab:seeds}
\end{table}

\noindent The three-condition gradient is preserved across all seeds.
\vota{}'s extremely low variance (std $= 0.15\pp$, range $= 0.3\pp$) reflects ceiling-proximate performance: only $\sim$10 of 1{,}190 samples are misclassified, and seed variation affects only which specific boundary cases fall on which side.
The low variance also suggests that \vota{}'s high information density creates a smooth loss landscape insensitive to initialization.

\end{document}
